\section{Цель работы и постановка задачи}
\label{sec:Chapter2} \index{Chapter2}

\subsection{Актуальность работы}
Рост скорости сетевого трафика и требование быстро менять поведение коммутаторов без замены оборудования делают актуальной программируемую сетевую обработку. Компиляция P4-программы в конвейер RMT включает этап размещения логических таблиц по стадиям при ограничениях памяти и графе зависимостей. Качество этого этапа влияет на число занятых стадий и задержку обработки пакета, поэтому исследование эвристик размещения и их реализации для типовых ограничений TDG сохраняет практический интерес для разработчиков компиляторов и сетевого оборудования.

\subsection{Цель работы}
Целью данной работы является исследование жадных алгоритмов решения задачи размещения логических таблиц классификации по стадиям конвейера RMT-коммутатора, а также разработка программной реализации этих алгоритмов для оценки ресурсных затрат P4-программ.

\subsection{Постановка задачи}
\begin{enumerate}
	\item Провести анализ существующих жадных эвристик (FFD, FFL, FFLS, MCF) и их модификаций.
	\item Разработать многопроходный алгоритм FFD с поддержкой горизонтального разрезания таблиц.
	\item Реализовать алгоритм FFLS (First Fit by Level and Size) как эталон для сравнения.
	\item Реализовать поддержку учёта всех четырёх типов зависимостей между таблицами (match, action, successor, reverse match).
	\item Провести экспериментальное исследование на синтетических конфигурациях, различающихся количеством таблиц и числом записей.
	\item Сравнить алгоритмы по метрикам времени выполнения и количества использованных стадий.
	\item Проанализировать полученные результаты и сформулировать рекомендации по выбору алгоритма в зависимости от характеристик входной программы.
\end{enumerate}

\subsection{Ограничения}
В работе рассматриваются жадные эвристики размещения без гарантии глобального оптимума числа стадий. Реализован прототип на языке Python~3 без привязки к промышленному компилятору конкретного чипа; модель аппаратных квот задаётся агрегированно (SRAM и TCAM на стадию). Эксперименты выполняются на синтетически генерируемых TDG с контролируемыми параметрами плотности и типов зависимостей, что ограничивает прямое перенесение численных выводов на реальные большие P4-программы без дополнительной калибровки.
